\documentclass[11pt,a4paper]{article}
\usepackage[T1]{fontenc}
\usepackage[utf8]{inputenc}
\usepackage[ngerman,english]{babel}
\usepackage{amsmath,amssymb,mathtools}
\usepackage{geometry}
\usepackage{enumitem}
\usepackage{hyperref}
\geometry{margin=2.5cm}
\setlength{\parindent}{0pt}
\setlength{\parskip}{0.6em}

\title{P09 --- Bost--Connes, Hochschild and Charged Cohomology}
\author{Objekt-X-Programm}
\date{9. August 2026}

\begin{document}
\maketitle

\begin{center}
\textbf{Status: SYN FINAL AUDITED}
\end{center}

\paragraph{Auditbasis.}
Pass-A: \texttt{AUDIT-2026-08-09-P09-PassA-FINAL-SEAL};
SYN-Primärcheck: \texttt{AUDIT-2026-08-09-P09-SYN-Primaercheck};
SYN-Zweitcheck: \texttt{AUDIT-2026-08-09-P09-SYN-Zweitcheck-Pfadgebunden}.
Der Scope ist der BC-/Hochschild-Strang NEU-174--219; NEU-222 wird ausschließlich als Superseding-Scan verwendet.

\begin{abstract}
P09 konsolidiert den Bost--Connes-/Hochschild-Strang des Objekt-X-Programms. Auf
$A_{\mathbb Q}^{\rm alg}$ existiert eine nichttriviale neutrale Hochschild-4-Klasse. Der geladene Fortschritt entsteht über eine singuläre faktoriale Potentialroute und eine korrigierte äußere Derivation
\[
 [D_g^{\rm corr}]\neq0\in HH^1(A_{\rm alg},A_{C^*})_g.
\]
Ein logarithmischer Koeffiziententyp und der globale Bimodul
$\mathfrak M_{\rm glob}^{\log}$ erlauben den nichttrivialen geladenen Cup
\[
 \boxed{[D_g^{\rm corr}]\smile[\Theta^\wedge]\neq0
 \in HH^4(A_{\rm alg},\mathfrak M_{\rm glob}^{\log})_g.}
\]
Die zyklische Verfeinerung ist nicht automatisch möglich. Für $\beta>1$ liefert eine gradneutralisierte KMS-Auswertung einen nichtverschwindenden getwisteten Hochschildkozykel; im adelischen Dilatationsmodell besitzt der kanonische Basislift dagegen keine globale konstante Rotationseigenrelation:
\[
 \boxed{t\Phi_0\neq C\Phi_0\qquad\forall C\in\mathbb C.}
\]
P09 konstruiert weder einen zyklischen Ersatzrepräsentanten noch eine Weil-/Gamma-, Gram-, Hilbert--Pólya- oder Objekt-X-Endstruktur und enthält keinen RH-Beweis.
\end{abstract}

\section{Typ-, Notations- und Statusfirewalls}

Im I1-Grundblock setzen wir
\[
 \boxed{A:=A_{\mathbb Q}^{\rm alg}}.
\]
In den späteren Dateien wird derselbe algebraische BC-Kern überwiegend mit $A_{\rm alg}$ bezeichnet. Drei Koeffiziententypen sind strikt zu unterscheiden:
\[
 A_{\rm alg}\subset A_{C^*},\qquad
 \mathcal A^{\log}\subset A_{C^*},\qquad
 \mathfrak M_{\rm glob}^{\log}\subset\mathcal A^{\log}.
\]
Eine nichttriviale Klasse mit Koeffizienten in $A_{C^*}$ oder $\mathfrak M_{\rm glob}^{\log}$ impliziert keine nichttriviale Selbstkoeffizientenklasse in $HH^*(A_{\rm alg},A_{\rm alg})$.

Historische Formeln mit $D_g(e(r))=0$ sind superseded. Bindend ist
\[
 \boxed{D_g^{\rm corr}(e(r))=\mu_m C_{m,n;r}\mu_n^*,\qquad g=m/n.}
\]
Die Grenzdefinition ist nur als punktweiser Normgrenzwert auf jedem festen $a\in A_{\rm alg}$ gebucht; gleichmäßige Konvergenz in Derivationsoperatornorm wird nicht behauptet.

Ebenso ist $MX_N\to0$ falsch. Korrekt ist
\[
 \boxed{M(0)=0\Longrightarrow MX_N\text{ ist für }N\gg1\text{ exakt konstant}.}
\]

Die Katalognummern NEU-191 und NEU-198 fehlen im Live-Bestand. Aus diesen fehlenden Dateien wird keine positive SYN-Aussage allein abgeleitet.

\section{Algebraischer BC/Hochschild-Grundblock}

Für vier verschiedene Primrichtungen liefert der neutrale Grundblock
\[
 \boxed{[\Omega_p]\neq0\in HH^4(A_{\mathbb Q}^{\rm alg},A_{\mathbb Q}^{\rm alg}),
 \qquad \deg_\Gamma\Omega_p=1_\Gamma.}
\]
Die Alternierung ist unnormalisiert; die ausgezeichnete Paarung besitzt den Faktor $4!=24$.

Daraus folgt nicht
\[
 HH^4(A_{\rm alg},A_{\rm alg})_g\neq0\qquad(g\neq1).
\]
Die geladene Klasse im separaten Vier-Prim-Polynommodell $S_p$ wird nicht als BC-Klasse migriert.

Der geladene Dualzyklus sieht den vollständig alternierenden Anteil $\operatorname{Alt}_4L$. Die frühe symmetrische Produktschablone besitzt $\operatorname{Alt}_4L=0$. Ein determinantischer Modellkandidat paart zwar mit Wert $24$, scheitert aber am Hochschildrand. Das ist ein Kandidaten-No-go, kein globaler No-go für geladenes $HH^4$.

\section{Singuläre Potentialroute und geladene HH1-Klasse}

Für jeden nichtneutralen reduzierten Grad $g=m/n$ gilt
\[
 \boxed{Z_g=\{0\}.}
\]
Die faktoriale Kette
\[
 L_j=(j+1)!,\qquad P_j=1_{L_j\widehat{\mathbb Z}}
\]
liefert $\operatorname{Sing}(X)=\{0\}$. Der ältere No-go gegen exakt unter allen Primtransporten geschlossene totale Teilbarkeitsketten widerspricht dieser Konstruktion nicht, weil die faktoriale Route mit einem normkontrollierten Transportband arbeitet.

Die korrigierte geladene Derivation erfüllt für $g\neq1$
\[
 \boxed{[D_g^{\rm corr}]\neq0\in HH^1(A_{\rm alg},A_{C^*})_g.}
\]
Ihre Nichtinnerheit wird durch einen Offdiagonaltest gegen beschränkte Implementierer nachgewiesen. Sie landet im Allgemeinen nicht in $A_{\rm alg}$; deshalb bleibt eine geladene Selbstkoeffizientenklasse in $HH^1(A_{\rm alg},A_{\rm alg})_g$ offen.

Normkonvergente Potentialimplementierer liefern nur innere beziehungsweise im relevanten Quotienten unsichtbare Derivationen. Ein globaler normstetiger $A_{\rm alg}$-Bimoduloperator
\[
 R:A_{C^*}\to\mathcal A^\infty\subsetneq A_{C^*}
\]
kann nicht als nichttrivialer universeller Glätter die Zieltypfrage lösen.

Die drei konkreten dyadischen L/R/S-Platzierungen aus NEU-205 bleiben Kandidaten-No-gos. Die relation-adaptierte, $N$-abhängige ``Architecture III'' bleibt ausdrücklich offen.

\section{Logarithmischer Koeffiziententyp und geladener HH4-Cup}

Die Reparatur des gescheiterten Schwartz-Zieltyps ist die direkte Konstruktion
\[
 \boxed{\mathcal B_{\rm alg}\subsetneq\mathcal B^{\log}\subsetneq C(\widehat{\mathbb Z})}
\]
mit stabilen Operationen $\sigma_k$, $\rho_k$, $T_a$ und $G_{a,d}\in\mathcal B^{\log}$. Daraus entsteht die graduierte algebraische $*$-Algebra $\mathcal A^{\log}$.

Es gilt
\[
 \boxed{D_g^{\rm corr}(A_{\rm alg})\subseteq\mathcal A^{\log},\qquad
 [D_g^{\rm corr}]\neq0\in HH^1(A_{\rm alg},\mathcal A^{\log})_g.}
\]
Für den globalen logarithmischen Bimodul folgt stärker
\[
 \boxed{[D_g^{\rm corr}]\neq0\in HH^1(A_{\rm alg},\mathfrak M_{\rm glob}^{\log})_g.}
\]
Lokale stärkere Behauptungen mit $M_{g,p}^{\log}$ werden ohne vollständige lokale Bimodultypisierung nicht migriert.

Der Mehrparameter-Følner-Nachweis konstruiert einen partiellen Modulquotienten, einen Dualzeugen und eine nichtverschwindende Paarung. Damit
\[
 \boxed{[D_g^{\rm corr}]\smile[\Theta^\wedge]\neq0
 \in HH^4(A_{\rm alg},\mathfrak M_{\rm glob}^{\log})_g.}
\]
Der Beweis benötigt nicht den vollen Quotienten
$\mathfrak M_{\rm glob}^{\log}/[A_{\rm alg},\mathfrak M_{\rm glob}^{\log}]$; dessen Struktur bleibt offen. Ebenso folgt keine Klasse in $HH^4(A_{\rm alg},A_{\rm alg})_g$.

\section{KMS-Neutralisierung und getwisteter Hochschildkozykel}

Für ein homogenes nichtneutrales Zielelement $\eta$ erzwingt die KMS-Gleichung
\[
 \boxed{\omega_\beta(\eta)=0\qquad(\beta>0).}
\]
Ein KMS-Zustand ist keine gewöhnliche Spur und annihiliert im Allgemeinen nicht $[A,M]$.

Nach expliziter Gradneutralisierung reduziert die ausgezeichnete Auswertung für $\beta>1$ auf
\[
 \omega_{\beta,\chi}(\sigma_P(G_q)),
\]
und für alle extremalen Gibbs-Zustände im bewiesenen Bereich gilt
\[
 \boxed{\omega_{\beta,\chi}(\sigma_P(G_q))>0.}
\]
Der Fall $\beta=1$ wird durch diese Gibbs-Rechnung nicht entschieden.

Für die Standard-Letztrandkonvention ist
\[
 \boxed{\sigma_\beta=\alpha_{-i\beta}=\theta_\beta^{-1}.}
\]
Im bewiesenen Gibbs-Bereich folgt
\[
 \boxed{b^{\sigma_\beta}\Phi_{\beta,\chi}=0,\qquad
 0\neq\Phi_{\beta,\chi}\in Z^4_{\sigma_\beta,\mathrm{Hoch}}(A_{\rm alg}).}
\]
Strukturell gilt für diese rohe I4-Kochain
\[
 \boxed{T_{\sigma_\beta}\Phi_{\beta,\chi}=g^{-\beta}\Phi_{\beta,\chi}.}
\]
Im bewiesenen Nichtnullbereich $\beta>1$ und für $g\neq1$ ist dieser konkrete Repräsentant deshalb nicht standardmäßig getwistet-zyklisch. Diese Aussage betrifft $\Phi_{\beta,\chi}$, nicht den späteren kanonischen Basislift $\Phi_0$.

\section{Parazyklische und Hopf-SAYD-No-gos}

Der $w$-Eigenraum von $T$ ist ein $b^\sigma$-Unterkomplex. Für $w\neq1$ wird er bei gewöhnlicher Invarianten-/Koinvarianten-Zyklisierung annihiliert, weil $1-T$ dort invertierbar ist.

Eine externe eindimensionale Eigenlinie kompensiert den $T$-Eigenwert nur formal und definiert keine zyklische Koeffiziententheorie. Für $\beta>0$ existiert außerdem kein eindimensionales unital-nichtdegeneriertes $\sigma_\beta$-äquivariantes $A_{\rm alg}$-Bimodul des benötigten Typs.

Die $\mathbb Q_+^\times$-Gradierung liefert kanonisch eine Koaktion, nicht bereits eine kanonische Gruppenalgebra-Aktion. Der minimale reparierte Hopf-Typ
\[
 \mathcal H_\beta=\mathbb C[\mathbb Z]
\]
wirkt durch $\sigma_\beta$. Im Standard-SAYD-Setup kollidieren exakter KMS-Twist, Ladungskompensation und Stabilität. Ein nichtstandardmäßiger $A$-relativer Hopf-Koeffizient bleibt offen.

\section{Adelische Dilatation, Morita-Ecke und Orbitmarkierung}

Die BC-Endomorphismen besitzen die adelische automorphe Dilatation
\[
 \widetilde B=C_0(\mathbb A_f),\qquad (\gamma_rF)(a)=F(r^{-1}a),
\]
mit
\[
 \widetilde A=C_0(\mathbb A_f)\rtimes_\gamma\mathbb Q_+^\times.
\]
Für $e=1_{\widehat{\mathbb Z}}$ gilt die Full-Corner-Realisierung. Für den gewählten algebraischen Kern ist zusätzlich explizit nachgerechnet
\[
 \boxed{e\widetilde A_{\rm alg}e=j_A(A_{\rm alg}).}
\]
Der Eckprojektor ist algebraisch voll, und der Kern besitzt lokale Einheiten.

Die unmarkierten gesättigten Orbitmodule kollabieren:
\[
 \boxed{N_k=N_0\qquad\forall k\in\mathbb Z.}
\]
Daher ist die globale unmarkierte Multiplikationsabbildung nicht injektiv. Die Orbitinformation wird extern markiert:
\[
 \mathcal N_{\rm tag}=\bigoplus_{k\in\mathbb Z}^{\rm alg}N_0\delta_k.
\]
Auf $\mathcal N_{\rm tag}$ existiert eine typkorrekte Eigenfamilie $\Omega_\lambda$. Die Multiplikation mit $U_{g^{-1}}$ erhält den Orbitindex und ist nicht der externe Shift $T^{-1}$.

\section{Kanonischer Basislift und endgültiger Rotations-No-go}

Der kanonische Ecklift ist
\[
 \boxed{\widetilde L_0=\eta_0\circ j_M\circ L^{\rm cup}:
 A_{\rm alg}^{\otimes4}\to I_0,}
\]
und
\[
 \boxed{\widetilde L_0\in Z^4(A_{\rm alg},I_0).}
\]
Die Recovery-Identität ist als Inklusion
\[
 \Pi_0\circ\eta_0=\iota_{M_0\hookrightarrow N_0}
\]
zu lesen; nach Eckkompression erhält man $\operatorname{id}_{M_0}$.

Der Lift trägt den BC-Grad $g$, lebt aber vollständig im Orbit-Nullsummanden:
\[
 \boxed{\kappa=0,\qquad\varepsilon=0.}
\]
Daher ist jedes Orbitgewicht $\lambda$ auf dem kanonischen Lift wirkungslos.

Die historischen Zwischenbehauptungen
\[
 t\Phi_0=g^{-\beta}\Phi_0,
 \qquad s=-1
\]
sind superseded. KMS, Twistkommutation und Gradinformation liefern keine Rotation von $L(a_0,a_1,a_2,a_3)$ zu $L(a_1,a_2,a_3,a_4)$.

Für das Unit-Slot-Tupel
\[
 (a_0,a_1,a_2,a_3,a_4)=(\mu_P^*,\mu_{p_1},\mu_{p_2},\mu_{p_3},1)
\]
gilt
\[
 \Phi_0(a_0,\ldots,a_4)=0,
\]
während für $\beta>1$
\[
 (t\Phi_0)(a_0,\ldots,a_4)
 =-\left(\prod_{i=1}^3\log p_i\right)n^{-\beta}\omega_{\beta,\chi}(G_P)\neq0.
\]
Folglich
\[
 \boxed{t\Phi_0\neq C\Phi_0\qquad\forall C\in\mathbb C.}
\]
Das schließt den kanonischen skalaren Basislift ab, nicht jedoch andere zyklische/getwistet-zyklische Repräsentanten, genuin orbitverschiebende nichtkanonische Lifte, andere Koeffizientenkategorien oder archimedische Korrekturen.

\section{Offene Endknoten}

Offen bleiben insbesondere:
\begin{enumerate}[label=\arabic*.]
\item eine geladene Selbstkoeffizientenklasse in $HH^1(A_{\rm alg},A_{\rm alg})_g$;
\item eine geladene Selbstkoeffizientenklasse in $HH^4(A_{\rm alg},A_{\rm alg})_g$;
\item die volle Struktur von $\mathfrak M_{\rm glob}^{\log}/[A_{\rm alg},\mathfrak M_{\rm glob}^{\log}]$;
\item die topologische Banach-/Fréchet-Vervollständigung des logarithmischen Zieltyps;
\item die lokalen Resttypisierungen [O-217-1d], [O-217-2b-5], [O-217-2c-5land];
\item der KMS-Grenzfall $\beta=1$ in der I4-Diagonalauswertung;
\item ein nichtstandardmäßiger $A$-relativer Hopf-Koeffizient;
\item [O-219-cyclic-representative];
\item ein genuin orbitverschiebender nichtkanonischer Lift;
\item [O-219-6], die Weil-/Gammafaktorpaarung;
\item die relation-adaptierte NEU-205-Architecture III.
\end{enumerate}

\section{Routing und Gesamturteil}

P10 darf isolierte Kandidaten-No-gos spiegeln; strukturtragende No-gos bleiben auch in P09. Die globale nichtorthogonale Gramkopplung, der Mediator zwischen Primkanälen und archimedischem Anteil und die Objekt-X-Gesamtgeometrie gehören nach P11. Der Weil-/Gamma-Pfad wird ab NEU-220 weitergeführt.

Der P09-Endstand lautet:
\[
 \boxed{\text{geladene Hochschildstruktur: positiv bis }HH^4
 \text{ mit Koeffizienten }\mathfrak M_{\rm glob}^{\log},}
\]
aber
\[
 \boxed{\text{kanonische skalare zyklische Rotation: negativ.}}
\]
Die Koeffizienten- und Zyklizitätsfirewalls verhindern einen automatischen Schluss auf eine zyklische Klasse, eine Weil-Identität oder einen Hilbert--Pólya-Operator.

\begin{center}
\fbox{\parbox{0.84\textwidth}{\centering
P09 enthält keinen RH-Beweis und konstruiert Objekt X nicht.}}
\end{center}

\end{document}
